\chapter{Strings and polygons}\label{ch:polygons}

\begin{goals}
\textbf{In this chapter} you open pit limits, lease boundaries, blast outlines and mapped contacts
as features with their attributes, from CSV string files, WKT, shapefiles and GeoJSON; copy them to
and from QGIS; draw and edit a polygon in plan; and export them back out.
\end{goals}

A \term{strings \& polygons} layer\index{polygons}\index{strings!with attributes} holds lines and
closed polygons in 3D, and its records are the \emph{features}: one pit limit, one lease, one blast.
Each feature carries its attributes, so the layer behaves like any other table in micro. Its
attribute table lists one row per feature (\key{t}), the filter keeps the features that match
(\cref{ch:filter}), the color list colors by any of its columns (\cref{ch:color}), and a click on a
line shows that feature's record (\cref{ch:picking}). A polygon has an inside: holes and multipart
features are kept, and the fill shows it.

\side{A DXF's design strings open as the segment-by-segment strings layer of \cref{sec:meshes-dxf};
\ui{Open as} reads the same drawing as features.}

\section{Opening}\label{sec:polygons-open}

Open the files like any other (\cref{ch:opening}). The open dialog names each one before it loads,
with its feature count: \ui{Shapefile · 12 polygons · 4 attributes · SIRGAS 2000 / UTM zone 23S →
strings \& polygons}.

\begin{description}
\item[A CSV string file] has an ID column (\col{STRING}, \col{ID}, \col{LINE}, \col{POLYGON},
  \col{FEATURE}…) and one row per vertex, in order, with \col{X}, \col{Y} and \col{Z}. The rows of
  one string sit together. This is how Leapfrog and Surpac export polylines. A string whose last
  vertex returns to its first is a polygon; any other columns become the feature's attributes, read
  from its first row. A table with an ID column whose values do not come in runs stays a block
  model or a table.
\item[A WKT-CSV table] has a column of well-known text (\col{wkt\_geom}, \col{WKT},
  \col{geometry}, \col{the\_geom}) beside its attributes. QGIS writes it when you save a layer as CSV
  with \code{GEOMETRY=AS\_WKT}, and its \menu{Edit>Copy Features} puts the same table on the
  clipboard. A \file{.wkt} file holds one geometry, or one per line.
\item[A shapefile] is several files: \file{.shp} (the shapes), \file{.dbf} (the attributes),
  \file{.prj} (the coordinate system) and often \file{.cpg} (the text encoding). Select them
  together; the open dialog makes them one layer. A \file{.zip} holding them opens directly.
\item[GeoJSON] opens a FeatureCollection, one Feature, or a bare geometry.
\item[A DXF] opens as features from \ui{Open as ▾} → \ui{strings \& polygons}: each polyline, line
  and circle is a feature, its DXF layer an attribute (\col{LAYER}).
\end{description}

Point geometries are not strings: a file mixing them skips the points and says how many, and a
file of only points opens as a points layer. A feature without a Z (a 2D shapefile, a WKT-CSV
digitized in plan) is drawn at $z=0$, and the status line says so; \menu{Geometry>Set elevation…}
places the layer at a bench.

\subsection{Longitude and latitude are refused}\index{CRS!geographic, refused}

micro places coordinates where they say and never reprojects (\cref{ch:opening}). A file in
degrees would land at $(-45, -20)$ on a metre grid, so it does not open. It is refused, with the
reason on the status line and on the open dialog's row, when its \file{.prj} names a geographic
system (\code{GEOGCS}), when a GeoJSON file has no projected \code{"crs"} (the GeoJSON standard is
longitude and latitude), when its SRID is a geographic EPSG code, or when every coordinate fits
within ±180 and ±90 over less than 20°. Reproject the file in your GIS (to UTM or your mine grid)
and open it again. Reprojection on open is planned.

\section{The layer}\label{sec:polygons-layer}

Lines draw as thin tubes in the feature's color, as design strings do. The right-click menu adds:

\begin{description}
\item[\menu{Display>Fill polygons}] shades each polygon's inside, holes left open, in its class
  color (or a single color when the layer is \ui{solid}). The fill follows the filter. It is shading,
  never picked: a click inside a polygon reaches what is behind.
\item[\menu{Display>Classes by…}] makes any text column the category, without reading the file
  again.
\item[\menu{Display>Line width…}] sets the tubes' half-width in metres.
\item[\menu{Geometry>Edit vertices}, \ui{Draw polygon into this layer}, \ui{Set elevation…}]
  are in \cref{sec:polygons-draw}.
\item[\menu{Data \& file>Export features…}, \ui{Copy features (WKT, for QGIS)}] are in
  \cref{sec:polygons-export}.
\end{description}

A filter such as \code{BENCH >= 850 and STATUS = "approved"} keeps whole features. The rectangle and
lasso select features (\cref{ch:picking}); the selection is what \ui{Copy features} and
\ui{Export features…} take when there is one. A vector section draws the features in its slab and,
on a plan, their fills (\cref{ch:composer}).

\section{QGIS and the clipboard}\label{sec:polygons-qgis}\index{QGIS!clipboard}

QGIS's clipboard is a table: a header row, tab-separated, the geometry as well-known text in a
\col{wkt\_geom} column first. micro reads and writes exactly that.

\begin{itemize}
\item \textbf{From QGIS to micro.} Select features in QGIS and press \key{Ctrl}+\key{C}. In micro,
  press \key{Ctrl}+\key{V} over the view (or \menu{File>Paste features (WKT from QGIS)}). The features
  arrive as a new layer, attributes as columns. Features digitized in 2D ask for the elevation to
  place them at; micro never guesses a Z.
\item \textbf{From micro to QGIS.} Right-click the layer, \menu{Data \& file>Copy features (WKT, for
  QGIS)}, and in QGIS \menu{Edit>Paste Features As>Temporary Scratch Layer}. The selected features go,
  or all of them when nothing is selected.
\end{itemize}

Coordinates cross unchanged in both directions, so both sides must already use the same projected
system.

\section{Drawing and editing}\label{sec:polygons-draw}\index{polygons!drawing}

\menu{Tools>Draw polygon…} and \menu{Tools>Draw line…} draw over whatever is open. Look from above
first (\key{d}, the plan view): each click lands where the view ray meets the drawing
plane, which is the plan section's elevation when a plan section is on, otherwise the top of the
scene.

\begin{itemize}
\item Click to add a vertex. A vertex near an existing feature vertex snaps to it, within ten pixels;
  hold \key{Alt} to place it freely.
\item Close a polygon by clicking its first vertex, double-clicking or pressing \key{Enter}.
  \key{Backspace} (or \key{Ctrl}+\key{Z}) takes the last vertex back; \key{Esc} stops.
\item micro asks for a name and, optionally, attributes as \code{key=value} pairs:
  \code{code=B12, bench=850}. Numbers stay numbers.
\end{itemize}

The first polygon starts a layer \ui{drawn polygons}; the next ones go into it while it is the active
layer. \ui{Draw polygon into this layer} on any strings \& polygons layer adds to that one instead,
with its columns.

\menu{Geometry>Edit vertices} puts a square on every vertex of the layer (of its selected features,
when there is a selection) and a small circle on every edge's midpoint. Drag a square to move a
vertex; its elevation stays. Drag a circle to add a vertex there. \key{Alt}-click a square, or click
it and press \key{Delete}, to remove it. Each edit is one step for \key{Ctrl}+\key{Z}. \key{Enter} or
\key{Esc} ends.

\begin{note}
A drawn or edited layer is kept in micro's own form, a WKT-CSV table. It lives in memory until you
save the project, which writes it into the project's \file{polygons/} folder (\cref{ch:projects});
the next open reads it back with its attributes. Editing a layer opened from a shapefile or GeoJSON
does not change that file: the layer becomes its WKT-CSV.
\end{note}

\begin{tryit}
Draw a limit over the demo model.
\begin{enumerate}
\item Open the demo project and press \key{d} for the plan view.
\item Choose \menu{Tools>Draw polygon…} and click five points around the high-grade core. Click the
  first point again to close.
\item Name it \code{pit limit}, with \code{phase=1} as its attribute.
\item Right-click \ui{drawn polygons} and choose \menu{Display>Fill polygons}.
\item Choose \menu{Geometry>Edit vertices}, drag one corner outward, then press \key{Ctrl}+\key{Z}
  to put it back and \key{Enter} to finish.
\end{enumerate}
\end{tryit}

\section{Exporting}\label{sec:polygons-export}

\menu{Data \& file>Export features…} writes the layer (or its selection) with every attribute:

\begin{description}
\item[WKT-CSV] (\file{.csv}): for QGIS, PostGIS and spreadsheets.
\item[GeoJSON] (\file{.geojson}): with a \code{"crs"} member, so the file is not read as longitude
  and latitude.
\item[Shapefile] (\file{.zip}): \file{.shp}, \file{.shx}, \file{.dbf}, \file{.cpg} and the
  \file{.prj} the layer came with. A shapefile holds one geometry type, so a layer with both
  polygons and lines writes two in the zip.
\item[DXF] (\file{.dxf}): one polyline per ring, on the DXF layer named by the category column, the
  attributes in each polyline's extended data under \code{MICRO}. micro reads them back, holes
  included.
\end{description}

All four open again in micro with the same features, attributes and coordinates
(\cref{sec:formats-polygons}).
